\section{Interpolations}

\subsection{Polynomial Interpolation}

\begin{exercise}
	Construct a Vandermonde matrix using the power basis and $n+1$ equidistant
	points in the interval $[-1,1]$, where $n=4,8,\ldots,40$.
	Specifically, the points are given by $x_i=-1+(i-1)h$, with $i=1,2,\ldots,n+1$
	and $h=2/n$. Use the Vandermonde matrix to determine the interpolating
	polynomials on these nodes for the functions $f_k(x)=\cos(kx)$ with
	$k=10,20,\ldots,100$. Study the value of these polynomials in $x=0$ for the
	different $n$ and $k$ values. What do you observe?
\end{exercise}

\begin{solution}

	The problem requires finding the coefficients $\mathbf{c} = (c_0, c_1, \dots,
		c_n)^T$ of the interpolating polynomial $P_n(x) = \sum_{j=0}^n c_j x^j$ of
	degree $n$. This is achieved by solving the linear system $V \mathbf{c} =
		\mathbf{f}$, where $V$ is the $(n+1) \times (n+1)$ Vandermonde matrix with
	entries $V_{i,j} = x_i^j$, and $\mathbf{f}$ is the vector of function values
	$f_k(x_i) = \cos(kx_i)$ evaluated at the equidistant nodes.

	To study the value of the polynomials at $x = 0$, we can exploit the
	definition of the power basis. Evaluating $P_n(x)$ at $x = 0$ leaves only the
	constant term:
	\[
		P_n(0) = \sum_{j=0}^n c_j 0^j = c_0
	\]
	Thus, studying the polynomial's value at the origin is equivalent to tracking
	the zeroth-order coefficient $c_0$. The absolute error at $x=0$ is computed as
	$\epsilon = |c_0 - 1|$, since the true analytical value is $\cos(k \cdot 0) = 1$
	for all $k$.

	\paragraph{Results}
	Three primary diagnostics were visualized to evaluate the interpolation
	scheme:
	\begin{enumerate}
		\item Condition Number: Figure \ref{fig:cond_number} demonstrates
		      the condition number $\kappa(V)$ as a function of the polynomial degree $n$. It
		      exhibits clear exponential growth.
		\item Interpolation Error at $x=0$: Figure \ref{fig:errors}
		      displays the base-10 logarithm of the absolute error $\log_{10}(|c_0 - 1|)$. A
		      feature in the data is a localized degradation of accuracy around $n=20$, where
		      the evaluation error spikes noticeably before settling or becoming entirely
		      consumed by numerical noise at higher degrees.
		\item Global Fit Quality: Figure \ref{fig:chi2} presents a heatmap
		      of the $\chi^2$ residuals. Consistent with the local error evaluation, the
		      global $\chi^2$ values exhibit a surge at $n=20$ across multiple wave numbers
		      $k$, indicating a severe structural breakdown of the global polynomial fit at
		      this specific degree.
	\end{enumerate}

	\begin{center}
		\centering
		\includegraphics[width=0.6\textwidth]{figures/3-1/condition-numbers.pdf}
		\captionof{figure}{Condition number of the Vandermonde matrix $\kappa(V)$ as
			a function of polynomial degree $n$, plotted on a logarithmic scale.}
		\label{fig:cond_number}
	\end{center}

	\begin{center}
		\centering
		\includegraphics[width=0.9\textwidth]{figures/3-1/errors-plots.pdf}
		\captionof{figure}{Left: Heatmap of the $\log_{10}$ absolute error at $x=0$
			for varying $n$ and $k$. Right: The computed value of $P_n(0)$ versus $n$ for
			different wave numbers $k$.}
		\label{fig:errors}
	\end{center}

	\begin{center}
		\centering
		\includegraphics[width=0.6\textwidth]{figures/3-1/chi-square.pdf}
		\captionof{figure}{Heatmap of the $\chi^2$ residuals representing the
			overall divergence of the interpolating polynomial from the true function
			$f_k(x) = \cos(kx)$.}
		\label{fig:chi2}
	\end{center}

	\paragraph{Discussion}
	The results highlight a complex interplay between sampling density, boundary
	effects, and numerical stability:
	\begin{itemize}
		\item Under-sampling (Small $n$, Large $k$): The highly oscillatory
		      nature of $f_k(x) = \cos(kx)$ requires a sufficient number of nodal points to
		      capture the waveform. When $n$ is too small relative to $k$, the equidistant
		      nodes fail to resolve the oscillations, leading to flat or highly inaccurate
		      approximations.
		\item The $n=20$ Transition Zone: The severe spike in both local
		      errors and global $\chi^2$ values at $n=20$ represents a critical numerical
		      transition. At this degree, the polynomial is flexible enough to trigger massive
		      oscillations at the boundaries, yet the nodal density is not high enough to lock
		      onto the tight frequencies of the larger $k$ modes.
		\item Ill-conditioning and Noise (Large $n$): As $n$ approaches 40,
		      $\kappa(V)$ becomes so astronomically large that the least-squares solver is
		      effectively fitting pure floating-point noise. While the theoretical error at
		      the center of the domain should be minimal, the catastrophic cancellation
		      inherent to the Vandermonde matrix entirely destroys the structural integrity of
		      the polynomial.
	\end{itemize}

	\paragraph{Conclusion}
	The exercise demonstrates that global polynomial interpolation on equidistant
	nodes using the standard power basis and Vandermonde matrices is fundamentally
	unstable for oscillatory functions. Increasing the polynomial degree to resolve
	higher frequencies does not converge to the true function; instead, it forces a
	transition through a highly volatile state (peaking around $n=20$) before
	eventually collapsing entirely due to catastrophic ill-conditioning at larger
	$n$.

\end{solution}

\subsection{Lagrange-Waring Interpolation}

\begin{exercise}
	In each case, interpolate the given function using \( n \) evenly spaced
	nodes in the given interval. Plot each interpolating function together with the
	exact function.
	\begin{enumerate}
		\item \( f(x) = \ln (x), \quad n = 2,3,4, \quad x\in [1,10] \)
		\item \( f(x) = \tanh (x), \quad n = 2,3,4, \quad x \in [-3,2] \)
		\item \( f(x) = \cosh (x), \quad n = 2,3,4, \quad x \in [-1,3] \)
		\item \( f(x) = |x|, \quad n = 3,5,7, \quad x \in [-2,1] \)
	\end{enumerate}
\end{exercise}

\begin{solution}

	To approximate the given functions, we implement polynomial interpolation over
	evenly spaced nodes. The interpolating polynomials are evaluated using the
	barycentric form of Lagrange interpolation.

	For a given interval $[a, b]$ and $n$ nodes, the evenly spaced interpolation
	nodes are generated as $x_k = a + k \frac{b-a}{n-1}$ for $k = 0, \dots, n-1$.
	The interpolation algorithm is evaluated over a high-resolution grid of 200
	points. To visualize the accuracy, the results are presented in a side-by-side
	format: the left panel overlays the exact function and the interpolant, while
	the right panel displays the absolute residual error, $E(x) = |P_{n-1}(x) -
		f(x)|$, on a logarithmic base-10 scale.

	\paragraph{Results}
	The resulting pseudo-$\chi^2$ error metrics for each function and node
	configuration are summarized in Table \ref{tab:interp_errors}. The visual
	representations of the interpolations and their corresponding absolute residuals
	are shown in Figures \ref{fig:interp_smooth} and \ref{fig:interp_abs}.

	\begin{center}
		\centering
		\captionof{table}{Pseudo-$\chi^2$ error metrics for polynomial interpolation
			using $n$ evenly spaced nodes.}
		\label{tab:interp_errors}
		\renewcommand{\arraystretch}{1.2}
		\begin{tabular}{l c r r r}
			\toprule
			Function          & Interval  & $n_1$ Error    & $n_2$ Error  & $n_3$ Error  \\
			\midrule
			$f(x) = \ln(x)$   & $[1, 10]$ & (n=2) 31.101   & (n=3) 3.285  & (n=4) 0.572  \\
			$f(x) = \tanh(x)$ & $[-3, 2]$ & (n=2) -29.363  & (n=3) 18.957 & (n=4) -3.529 \\
			$f(x) = \cosh(x)$ & $[-1, 3]$ & (n=2) 1171.471 & (n=3) 26.629 & (n=4) 2.105  \\
			$f(x) = |x|$      & $[-2, 1]$ & (n=3) 112.930  & (n=5) 30.394 & (n=7) 5.374  \\
			\bottomrule
		\end{tabular}
	\end{center}

	\begin{center}
		\centering
		\includegraphics[width=0.8\textwidth]{figures/3-2/log(x)-plot.pdf}
		\vspace{0.5cm} \\
		\includegraphics[width=0.8\textwidth]{figures/3-2/tanh(x)-plot.pdf}
		\vspace{0.5cm} \\
		\includegraphics[width=0.8\textwidth]{figures/3-2/cosh(x)-plot.pdf}
		\captionof{figure}{Side-by-side evaluation of polynomial interpolations and
			absolute residual errors for the smooth functions $\ln(x)$, $\tanh(x)$, and
			$\cosh(x)$.}
		\label{fig:interp_smooth}
	\end{center}

	\begin{center}
		\centering
		\includegraphics[width=0.8\textwidth]{figures/3-2/abs_x-plot.pdf}
		\captionof{figure}{Interpolation and absolute residual error for the
			non-smooth function $f(x) = |x|$. Note the localized error peak near the
			non-differentiable cusp at $x=0$.}
		\label{fig:interp_abs}
	\end{center}

	\paragraph{Conclusion}
	The side-by-side visualizations effectively confirm that barycentric Lagrange
	interpolation successfully captures the global behavior of the evaluated
	functions for low-degree polynomials. However, the residual analysis is far more
	reliable than the global pseudo-$\chi^2$ metric. As seen with $\tanh(x)$, the
	$\chi^2$ metric yields non-physical negative values because the exact function
	crosses zero. The logarithmic residual plots provide a much clearer, strictly
	positive evaluation of the spatial distribution of the interpolation error.

\end{solution}

\subsection{Error on Interpolation}

\begin{exercise}
	\begin{enumerate}
		\item For each case below, compute the polynomial interpolant using $n$
		      second-kind Chebyshev nodes in $[-1,1]$ for $n=4,8,12,\ldots,60$. At each value
		      of $n$, compute the $\infty$-norm error (that is,
		      $||f-p||_\infty=\max_{x\in[-1,1]} |p(x)-f(x)|$ evaluated for at least 4000
		      values of $x$). Using a log-linear scale, plot the error as a function of $n$,
		      then determine a good approximation to the constant $K$ in {eq}`eq-spectral`.
		      \begin{enumerate}
			      \item $f(x) = 1/(25x^2+1)$
			      \item $f(x) = \tanh(5 x+2)$
			      \item $f(x) = \cosh(\sin x)$
			      \item $f(x) = \sin(\cosh x)$
		      \end{enumerate}
		\item Make an analogous plot using $n$ equidistant points for the
		      interpolation.
	\end{enumerate}
\end{exercise}

\begin{solution}

	To evaluate the impact of node distribution on error propagation, we contrast
	equidistant (Lagrange) nodes with Chebyshev nodes.

	The approximation quality is measured using the $L^\infty$-norm,
	computationally approximated by evaluating the absolute error over a dense
	uniform grid of $4000$ evaluation points. Note that calculating the absolute
	error rather than the relative error ensures our convergence metrics are not
	skewed by the roots of the target function, which is mathematically robust when
	evaluating the global uniform norm. The theoretical spectral error bound behaves
	as:
	\[
		\|f - p_n\|_\infty \sim \mathcal{O}(K^{-n})
	\]
	Taking the natural logarithm yields $\ln \|f - p_n\|_\infty \approx -n \ln K$.
	Thus, the constant $K$ can be estimated as $K = \exp(-m)$, where $m$ is the
	slope extracted from a linear fit of the bounding error data on a log-linear
	scale.

	\paragraph{Results}
	The estimated spectral constants $K$ alongside their corresponding logarithmic
	decay slopes $m$ are tabulated in Table~\ref{tab:interp_errorss}. Following
	this, Figures~\ref{fig:runge_cheb} through \ref{fig:sincosh_lag} visually map
	the interpolating curves and the $L^\infty$-norm progression as a function of
	the polynomial degree $n$.

	\begin{center}
		\centering
		\captionof{table}{Estimated spectral convergence constant $K$ and
			logarithmic slope $m$ computed over the interval $n \in[4, 60]$.}
		\label{tab:interp_errorss}
		\renewcommand{\arraystretch}{1.2}
		\begin{tabular}{lcccc}
			\toprule
			Function        & \multicolumn{2}{c}{Chebyshev Nodes} & \multicolumn{2}{c}{Equidistant Nodes}                        \\
			\cmidrule(lr){2-3} \cmidrule(lr){4-5}
			                & $K$                                 & Slope ($m$)                           & $K$    & Slope ($m$) \\
			\midrule
			$1/(25x^2 + 1)$ & 1.2035                              & $-0.1852$                             & 0.7280 & 0.3175      \\
			$\tanh(5x + 2)$ & 1.2020                              & $-0.1840$                             & 0.8951 & 0.1108      \\
			$\cosh(\sin x)$ & 1.1513                              & $-0.1409$                             & 0.9728 & 0.0276      \\
			$\sin(\cosh x)$ & 1.1722                              & $-0.1589$                             & 0.9845 & 0.0156      \\
			\bottomrule
		\end{tabular}
	\end{center}

	\begin{center}
		\centering
		\begin{minipage}{0.48\textwidth}
			\centering
			\includegraphics[width=\linewidth]{figures/3-4/fa-Barycentric Chebyshev Interpolation.pdf}
			\captionof{figure}{Chebyshev interpolation: $1/(25x^2+1)$.}
			\label{fig:runge_cheb}
		\end{minipage}\hfill
		\begin{minipage}{0.48\textwidth}
			\centering
			\includegraphics[width=\linewidth]{figures/3-4/fa-Barycentric Lagrange Interpolation.pdf}
			\captionof{figure}{Equidistant interpolation: $1/(25x^2+1)$.}
			\label{fig:runge_lag}
		\end{minipage}
	\end{center}

	\begin{center}
		\centering
		\begin{minipage}{0.48\textwidth}
			\centering
			\includegraphics[width=\linewidth]{figures/3-4/fb-Barycentric Chebyshev Interpolation.pdf}
			\captionof{figure}{Chebyshev interpolation: $\tanh(5x+2)$.}
		\end{minipage}\hfill
		\begin{minipage}{0.48\textwidth}
			\centering
			\includegraphics[width=\linewidth]{figures/3-4/fb-Barycentric Lagrange Interpolation.pdf}
			\captionof{figure}{Equidistant interpolation: $\tanh(5x+2)$.}
		\end{minipage}
	\end{center}

	\begin{center}
		\centering
		\begin{minipage}{0.48\textwidth}
			\centering
			\includegraphics[width=\linewidth]{figures/3-4/fc-Barycentric Chebyshev Interpolation.pdf}
			\captionof{figure}{Chebyshev interpolation: $\cosh(\sin x)$.}
		\end{minipage}\hfill
		\begin{minipage}{0.48\textwidth}
			\centering
			\includegraphics[width=\linewidth]{figures/3-4/fc-Barycentric Lagrange Interpolation.pdf}
			\captionof{figure}{Equidistant interpolation: $\cosh(\sin x)$.}
		\end{minipage}
	\end{center}

	\begin{center}
		\centering
		\begin{minipage}{0.48\textwidth}
			\centering
			\includegraphics[width=\linewidth]{figures/3-4/fd-Barycentric Chebyshev Interpolation.pdf}
			\captionof{figure}{Chebyshev interpolation: $\sin(\cosh x)$.}
		\end{minipage}\hfill
		\begin{minipage}{0.48\textwidth}
			\centering
			\includegraphics[width=\linewidth]{figures/3-4/fd-Barycentric Lagrange Interpolation.pdf}
			\captionof{figure}{Equidistant interpolation: $\sin(\cosh x)$.}
			\label{fig:sincosh_lag}
		\end{minipage}
	\end{center}

	\paragraph{Discussion}
	\begin{itemize}
		\item Equidistant Nodes \& Runge's Phenomenon: The uniformly spaced
		      nodes lead to highly erratic polynomial oscillations near the boundaries $x=\pm
			      1$. This catastrophic cancellation phenomenon is strictly captured by the data:
		      $K < 1$ and $m > 0$ for all models, highlighting that the $L^\infty$-norm
		      diverges exponentially as the polynomial degree $n$ increases.

		\item Chebyshev Nodes \& Spectral Convergence: Deploying Chebyshev
		      points effectively mitigates the boundary oscillations. The values of $K > 1$
		      reflect the expected theoretical bounds for exponential error decay. The Runge
		      function $1/(25x^2+1)$ holds complex poles near the real axis. By contrast,
		      entire mathematical functions devoid of poles, such as $\cosh(\sin x)$,
		      theoretically guarantee super-exponential convergence.
	\end{itemize}

	\paragraph{Conclusion}
	The analysis effectively bridges numerical experiments with complex analysis
	theory, demonstrating the immense significance of interpolation grid
	conditioning. Uniform spacing predictably incurs Runge's phenomenon and
	exponential divergence for large degrees. In stark contrast, optimally
	clustering interpolation nodes around interval boundaries using a Chebyshev
	distribution ensures mathematical stability and recovers powerful spectral
	accuracy.

\end{solution}

\begin{exercise}
	The Chebyshev points can be used when the interval of interpolation is
	$[a,b]$ rather than $[-1,1]$ by means of the change of variable
	\begin{equation*}
		z = \psi(x) = a + (b-a)\frac{(x+1)}{2}.
	\end{equation*}
	Plot a polynomial interpolant of $f(z) =\cosh(\sin z)$ over $[0,2\pi]$ using
	$n=40$ Chebyshev nodes.
\end{exercise}

\begin{solution}

	To interpolate the function $f(x) = \cosh(\sin x)$ over the arbitrary interval
	$[a, b] = [0, 2\pi]$, we first generate the standard Chebyshev nodes of the
	first kind on the canonical interval $[-1, 1]$:
	\[
		t_i = \cos\left(\frac{2i - 1}{2n}\pi\right), \quad i = 1, \dots, n
	\]
	These nodes are subsequently mapped to the target interval $[a, b]$ via the
	affine transformation:
	\[
		x_i = \frac{b - a}{2} t_i + \frac{a + b}{2}
	\]
	The polynomial interpolant is evaluated using the barycentric formulation of
	Lagrange interpolation. When implemented correctly with Chebyshev weights, this
	provides a numerically stable algorithm capable of evaluating the interpolant in
	$\mathcal{O}(n)$ operations.

	\paragraph{Results}
	The interpolation was performed using $n=40$ nodes over 1000 evaluation
	points. The maximum absolute error over the interval was found to be:
	\[
		\max_{x \in [0, 2\pi]} |f(x) - p_{40}(x)| \approx 5.468 \times 10^{-5}
	\]
	The resulting interpolant, alongside the target function evaluated on a
	logarithmic scale, is visualized in Figure \ref{fig:cheb_interp2}.

	\begin{center}
		\centering
		\includegraphics[width=0.8\textwidth]{figures/3-4/cosh(sinh(x)) graph.pdf}
		\captionof{figure}{Barycentric Chebyshev interpolation of $f(x) = \cosh(\sin
				x)$ over $[0, 2\pi]$ using $n=40$ nodes. The function $y$ is plotted on a
			logarithmic scale.}
		\label{fig:cheb_interp2}
	\end{center}

	\paragraph{Discussion and Conclusion}
	The results confirm the absence of the Runge phenomenon, successfully
	demonstrating that clustering nodes near the domain boundaries mitigates the
	unbounded oscillations typical of equispaced Lagrange interpolation. The
	observed maximum error of roughly $5.5 \times 10^{-5}$ for $n=40$ is consistent
	with the expected behavior of polynomial interpolation for this function. In
	particular, although $f(x) = \cosh(\sin x)$ is entire, its rapid growth in the
	complex plane limits the convergence rate of the interpolant. Consequently,
	achieving errors close to machine precision at moderate polynomial degrees is
	not expected. Moreover, as the number of nodes increases, the error does not
	necessarily decrease monotonically; beyond a certain point, numerical effects
	such as rounding errors and loss of significance can lead to a slight increase
	in the observed maximum error. Overall, the numerical results are in good
	agreement with theoretical expectations.

\end{solution}

\begin{exercise}

	Let $x_1,\ldots,x_n$ be standard Chebyshev points. These map to the $z$
	variable as $z_i=\phi(x_i)$ for all $i$, where $\phi$ is the transformation
	defined in
	\[
		z = \phi(x) = \frac{2x}{1-x^2}.
	\]
	Suppose that $f(z)$ is a given function whose domain is the entire real line.
	Then the function values $y_i=f(z_i)$ can be associated with the Chebyshev nodes
	$x_i$, leading to a polynomial interpolant $p(x)$. This in turn implies an
	interpolating function on the real line, defined as
	\[
		q(z)=p\bigl(\phi^{-1}(z)\bigr) = p(x)\,.
	\]
	Implement this idea to plot an interpolant of $f(z)=(z^2-2z+2)^{-1}$ using
	$n=30$. Your plot should show $q(z)$ evaluated at 1000 evenly spaced points in
	$[-6,6]$, with markers at the nodal values (those lying within the $[-6,6]$
	window). [Hint: If you prefer avoid dealing with potential infinities consider
			using the Chebyshev nodes of the first kind.]

\end{exercise}

\begin{solution}

	To interpolate $f(z)$ on the infinite domain, we utilize a mapping $\phi: (-1,
		1) \to \mathbb{R}$. We employ Chebyshev nodes of the first kind to avoid the
	singularities at $x = \pm 1$. The inverse mapping is defined as:
	\[
		x = \phi^{-1}(z) = \frac{-1 + \sqrt{1 + z^2}}{z}, \quad z \neq 0
	\]
	The interpolation $p(x)$ is computed using the barycentric form of the
	Lagrange interpolant, which provides superior numerical stability for high-order
	polynomials ($n=30$).

	\paragraph{Results and Visualization}
	The maximum absolute error recorded over the interval $z \in [-6, 6]$ was
	$\|f(z) - q(z)\|_{\infty} \approx 4.399 \times 10^{-4}$. As illustrated in
	Figure \ref{fig:cheb_interp}, the interpolant $q(z)$ tracks the target function
	with high fidelity.

	\begin{center}
		\centering
		\includegraphics[width=0.85\textwidth]{figures/3-4/cheb_interpolation_plot.pdf}
		\captionof{figure}{Interpolation of $f(z) = (z^2-2z+2)^{-1}$ on the real
			line using $n=30$ transformed Chebyshev nodes. The markers indicate the nodal
			points $z_i$ that fall within the viewing window.}
		\label{fig:cheb_interp}
	\end{center}

	\paragraph{Discussion}
	The high accuracy achieved with only 30 nodes is a consequence of the
	analyticity of $f(z)$ in a neighborhood of the real axis. The transformation
	$\phi(x)$ effectively concentrates nodes near the "feature" of the function (the
	peak near $z=1$), while naturally sparsifying nodes in the tails where the
	function decays towards zero. This prevents the numerical instabilities
	typically encountered when using large numbers of equidistant points on wide
	domains.

\end{solution}

\begin{exercise}

	Spectral theorem assumes that the function being approximated has infinitely
	many derivatives over $[-1,1]$. But now consider the family of functions
	$f_m(x)=|x|^m$.
	\begin{enumerate}
		\item How many continuous derivatives over $[-1,1]$ does $f_m$ possess?
		\item Compute the polynomial interpolant using $n$ second-kind Chebyshev
		      nodes in $[-1,1]$ for $n=10,20,30,\ldots,100$. At each value of $n$, compute the
		      max-norm error (that is, $\max |p(x)-f_m(x)|$ evaluated for at least 4000 values
		      of $x$). Using a single log-log graph, plot the error as a function of $n$ for
		      all six values $m=1,3,5,7,9,11$.
		\item Based on the results of parts (a) and (b), form a hypothesis about the
		      asymptotic behavior of the error for fixed $m$ as $n\rightarrow \infty$.
	\end{enumerate}

\end{exercise}

\begin{solution}

	To evaluate the regularity of the function family $f_m(x) = |x|^m$ over the
	interval $[-1, 1]$ for odd integers $m$, we consider its derivatives. For any
	given odd $m$, the function is continuously differentiable up to order $m-1$.
	The $(m-1)$-th derivative is proportional to $|x|$, which is strictly continuous
	across the domain. However, the $m$-th derivative is proportional to
	$\operatorname{sgn}(x)$, which exhibits a jump discontinuity at the origin
	$x=0$. Therefore, mathematically, $f_m \in C^{m-1}([-1,1])$; it possesses
	exactly $m-1$ continuous derivatives.

	\paragraph{Numerical Methodology}
	The loss of infinite differentiability implies that we can no longer expect
	spectral convergence. To quantify the convergence numerically, we approximate
	$f_m(x)$ using barycentric polynomial interpolation over $n$ second-kind
	Chebyshev nodes, defined algebraically as:
	\[
		x_j = -\cos\left(\frac{j-1}{n-1}\pi\right), \quad j=1,\dots,n
	\]
	The $L^\infty$-norm (maximum absolute error) is approximated by evaluating the
	interpolant over a dense, uniform grid of $4000$ points. The relationship
	between the error and the degree $n$ is evaluated on a log-log scale. For
	algebraic convergence of the form $\|f_m - p_n\|_\infty \approx C n^{-\nu}$,
	taking the natural logarithm yields:
	\[
		\ln \|f_m - p_n\|_\infty \approx -\nu \ln n + \ln C
	\]
	The convergence rate $\nu$ corresponds to the negative slope of the plot.

	\paragraph{Results}
	The computed slopes from the numerical experiments are compared against
	theoretical predictions in Table~\ref{tab:convergence_rates}.
	Figure~\ref{fig:convergence_rates} visualizes the algebraic convergence
	corresponding to the log-log linear decay.

	\begin{center}
		\centering
		\captionof{table}{Observed versus predicted convergence rates (log-log
			slopes) for $f_m(x)=|x|^m$.}
		\label{tab:convergence_rates}
		\renewcommand{\arraystretch}{1.2}
		\begin{tabular}{rccc}
			\toprule
			$m$ & Continuous Derivatives ($m-1$) & Observed Slope & Predicted Slope \\
			\midrule
			1   & 0                              & $-1.051$       & $-1$            \\
			3   & 2                              & $-3.147$       & $-3$            \\
			5   & 4                              & $-5.324$       & $-5$            \\
			7   & 6                              & $-7.660$       & $-7$            \\
			9   & 8                              & $-10.418$      & $-9$            \\
			11  & 10                             & $-12.724$      & $-11$           \\
			\bottomrule
		\end{tabular}
	\end{center}

	\begin{center}
		\centering
		\includegraphics[width=0.85\linewidth]{figures/3-4/convergence-rates.pdf}
		\captionof{figure}{Log-log plot demonstrating algebraic convergence of
			Chebyshev polynomial interpolation for functions with limited regularity. The
			steepness of the slope corresponds strictly to the parameter $m$.}
		\label{fig:convergence_rates}
	\end{center}

	\paragraph{Discussion \& Asymptotic Hypothesis}
	Spectral accuracy relies on a function being analytic or, minimally,
	infinitely differentiable. In contrast, $f_m(x) = |x|^m$ inherently contains a
	geometric "kink" at the origin $x=0$, fundamentally restricting its smoothness.

	As predicted by classical approximation theory, the rate at which polynomial
	interpolation converges is intimately linked to the function's regularity. Based
	on our numerical findings—where the observed log-log slope strongly correlates
	with $-m$—we formulate the following asymptotic hypothesis:
	\[
		\lim_{n \to \infty} \|f_m - p_n\|_\infty \sim \mathcal{O}(n^{-m})
	\]
	This algebraic convergence dictates that interpolating functions with local
	singularities (or discontinuous higher-order derivatives) inherently bottlenecks
	numerical accuracy to a polynomial rate, explicitly circumventing the
	exponential decay seen in pure spectral methods.

\end{solution}

\subsection{Some observations about the algorithms used in this section}

The computational efficiency and numerical stability of interpolation methods
vary significantly depending on the chosen basis and node distribution.

\paragraph{Direct Linear Systems}
Solving the polynomial interpolation problem by setting up a general $n \times
	n$ linear system, such as using the power basis with a Vandermonde
	matrix $\mathbf{V}_n \mathbf{c} = \mathbf{f}$, is often impractical for large
$n$. The cost of solving such a system using general-purpose decompositions like
LU or QR scales as $\propto n^3$.

It is crucial to distinguish between interpolation ($m=n$) and
least squares ($m > n$). While QR decomposition is slightly more
computationally demanding than LU, it is highly recommended for least squares
problems because it is far more stable than solving normal equations. For an $m
	\times n$ matrix, the thin QR decomposition requires approximately $2mn^2$
floating-point operations (flops).

\paragraph{Lagrange and Barycentric Forms}
Standard Lagrange interpolation is frequently cited in textbooks for its
theoretical elegance but criticized for its computational inefficiency. In its
basic form, evaluating the interpolating polynomial $p(x)$ at a single point
requires recomputing $n$ Lagrange basis functions $\ell_j(x)$, each involving a
product of $n-1$ terms. This leads to an evaluation cost of $\propto n^2$.

The barycentric form of Lagrange interpolation overcomes these
limitations by separating the calculation into a pre-computation step and an
evaluation step:
\begin{equation*}
	p(x) = \Phi_{n}(x) \sum_{j=1}^{n} \frac{\lambda_{j}}{x - x_{j}}f(x_{j})
\end{equation*}
The support coefficients (or weights) $\lambda_j$ depend only on the node
positions. Computing all $n$ weights requires $O(n^2)$ operations, but once they
are known, evaluating $p(x)$ for any $x$ costs only $\propto n$ flops.
Furthermore, the barycentric form is exceptionally efficient for
updating the interpolant; adding a new point $x_{k}$ costs only
$\propto (k-1)$ flops to update the weights.

\paragraph{The Chebyshev Advantage}
For certain node distributions, the computational burden is even lower. For
equidistant or Chebyshev points, analytic formulas allow the
weights $\lambda_j$ to be calculated in only $O(n)$ operations. In these cases,
the barycentric formula is considered superior to almost all other interpolation
schemes.

Chebyshev points, in particular, provide several critical advantages:
\begin{itemize}
	\item Improved Conditioning: Chebyshev nodes significantly improve
	      the conditioning of the interpolation problem compared to equidistant nodes,
	      making the result far less sensitive to rounding errors.
	\item Numerical Stability: When paired with the barycentric form,
	      Chebyshev nodes of the second kind demonstrate excellent stability. Potential
	      numerical issues near nodes are mitigated as errors in the numerator and
	      denominator essentially cancel out.
	\item Avoidance of Precision Loss: Unlike equidistant nodes, which
	      lead to Runge's phenomenon and a massive loss of accurate digits as $n$
	      increases, Chebyshev points distribute nodes in a way that minimizes error
	      amplification and maintains structural integrity even for high degrees.
\end{itemize}



